\documentclass{article}
\usepackage[T1]{fontenc}
\usepackage{lmodern}
\usepackage{graphicx}
\usepackage{amsmath,amssymb}
\usepackage[a4paper,margin=2cm]{geometry}
\usepackage[absolute,overlay]{textpos}
\setlength{\TPHorizModule}{1mm}
\setlength{\TPVertModule}{1mm}
\usepackage[indonesian]{babel}
\usepackage{hyperref}
\setlength{\emergencystretch}{2em}
% unit: Halaman 41

\begin{document}

% === Halaman 41 (python/native) ===

Elliptic Curve Diffie-Hellman (ECDH)

• Public: Kurva eliptik dan titik B(x,y) pada kurva

• Secret: Integer milik Alice, a, dan integer milik Bob, b

Alice, A

Bob, B

aB

bB

• Alice menghitung K = a (b.B) 

• Bob menghitung K = b(aB)

• Hasil perhitungan akan sama karena ab = ba

*) Sumber bahan: Debdeep Mukhopadhyay, Elliptic Curve Cryptography ,

 Dept of Computer Sc and Engg IIT Madras

41

% deskripsi: gambar native xref 162
\noindent\includegraphics[width=0.85\textwidth]{../../images/python/page_041_img_01.png}

% deskripsi: gambar native xref 163
\noindent\includegraphics[width=0.85\textwidth]{../../images/python/page_041_img_02.png}

\end{document}
